\section{Introduction}
\label{sec:introduction}

\IEEEPARstart{H}{andwritten} Text Recognition (HTR) has matured to the point where character error rates (CER) on the standard IAM benchmark~\cite{marti2002iam} sit in the low single digits: the CNN--RNN--CTC baseline of Retsinas~\etal~\cite{retsinas2022htrbestpractices} reaches ${\sim}4.6\%$~CER without any pretraining or synthetic data. With recognition largely solved on the widely used benchmarks, the open question in HTR shifts from \emph{whether} the model reads the line correctly to \emph{what the model looks at when it reads}, and whether those internal signals are clean enough to reuse for other document-analysis tasks such as writer identification, style analysis, and palaeographic dating that traditionally required hand-designed feature pipelines~\cite{souibgui2022vlac}.

This report follows the second question. The specific brief from the supervisor, restated verbatim from the project's opening emails, was: \emph{``make HTR explainable as much as possible for potential downstream tasks such as writer identification --- attention maps should be fine-grained and meaningful, and register tokens from the Vision Transformers Need Registers paper should be integrated to test their effect on attention quality.''} The rest of Section~\ref{sec:introduction} unpacks that brief into a background, a research gap, four research questions, and a list of contributions.

\subsection{Background and Motivation}
\label{sec:background}

Modern HTR pipelines fall into two structural families. The first is the CNN--BiLSTM--CTC family, canonically represented by Retsinas~\etal~\cite{retsinas2022htrbestpractices}: a convolutional stem collapses the input line image to a one-dimensional sequence, a recurrent head models sequential dependencies, and Connectionist Temporal Classification (CTC)~\cite{graves2006ctc} handles the length mismatch between output characters and input positions. This family is small, fast, and works well on limited training data because CNN inductive biases (locality, translation equivariance) match the physical structure of ink strokes. The second family is Transformer-based: either encoder-only with CTC (as in HTR-VT~\cite{li2024htrvt}) or encoder--decoder with cross-attention (as in TrOCR~\cite{li2023trocr}). These trade CNN inductive bias for a global self-attention mechanism that, in principle, exposes exactly the kind of internal signal --- \emph{where on the image is each output character being read from} --- that downstream applications need.

In practice, this internal signal is not clean. Darcet~\etal~\cite{darcet2024registers} showed that Vision Transformers trained on standard classification tasks develop a small number of anomalously high-norm patch tokens that behave as \emph{attention sinks}: they absorb a disproportionate share of the attention mass across layers, produce visible artifacts in attention maps, and hurt any downstream task that reads those maps. Their proposed fix is minimally invasive --- prepend a small number of learnable \emph{register tokens} to the patch sequence, and the sinks migrate onto the registers, leaving the patches free to produce interpretable attention. On classification ViTs the fix is clean: it improves downstream detection and segmentation while leaving classification accuracy essentially unchanged.

Whether the same fix works for a CTC-based HTR system is not obvious, and this is the empirical question the report addresses. HTR self-attention differs from classification self-attention in two structural ways. First, HTR has no dedicated CLS query: the queries at inference are the patch tokens themselves, so the ``sink'' phenomenon --- if it appears at all --- is queried by the same population of tokens that produce it, which is a different setting than the CLS-vs-patch dichotomy Darcet~\etal\ analyse. Second, the CTC loss is agnostic to attention structure: it only cares whether some monotonic alignment of the character labels through the token sequence has high probability, and does not directly penalize an attention map that ``looks bad'' as long as the decoding is correct. Both differences are reasons to expect that porting registers from classification ViTs to CTC HTR is not free.

\subsection{Problem Statement}
\label{sec:problem_statement}

The project has one deliverable and one measurable question.

\begin{itemize}
    \item \textbf{Deliverable.} A modular, reproducible HTR training and evaluation pipeline that supports CNN--RNN, ViT (from scratch), and pretrained-backbone models within a single interface, with attention-extraction hooks aligned to the visualization style of Fig.~5 in \emph{Beyond Memorization}~\cite{nikolaidou2024beyondmem}.
    \item \textbf{Measurable question.} Does adding $R \in \{0, 2, 4, 8, 16\}$ register tokens to a CTC ViT HTR system (a)~change the character/word error rate on the IAM Aachen split and the READ2016 test set, and (b)~change the quality of the attention maps produced by the encoder, judged both qualitatively (per-character maps in the style of \cite{nikolaidou2024beyondmem}) and quantitatively (attention concentration on the CTC decoding column of each character)?
\end{itemize}

\subsection{Research Gap}
\label{sec:research_gap}

The registers-for-attention idea is well established for classification ViTs and for dense-prediction tasks derived from them (segmentation, detection)~\cite{darcet2024registers}. It has not been systematically evaluated for HTR, which uses a sequence-labelling loss (CTC) rather than a per-image class posterior. Two consequences follow. First, no prior work quantifies the recognition-accuracy cost or benefit of registers in the CTC HTR setting. Second, no prior work characterizes what registers actually \emph{do} to attention in that setting --- whether the sink pattern documented for classification ViTs appears at all in CTC ViTs, and if so whether registers absorb it as cleanly as they do for classification.

The nearest concurrent work, HTR-VT~\cite{li2024htrvt}, is a CTC ViT for HTR and validates that ViT encoders can match CNN baselines on modern benchmarks, but it does not use registers and does not analyse attention quality. TrOCR~\cite{li2023trocr} exposes clean cross-attention maps because it is encoder--decoder, but that architectural choice is orthogonal to the registers question and comes with a $10\times$ larger parameter budget. The literature therefore leaves a gap that the CTC ViT + registers combination fills naturally: a small, fast HTR encoder whose self-attention is expected to be interpretable by construction.

\subsection{Research Questions}
\label{sec:research_questions}

The report answers four questions in order:

\begin{enumerate}
    \item \textbf{RQ1 (Recognition).} How does the register count $R$ affect CER and WER on IAM and READ2016, and are any observed differences larger than the seed-noise floor?
    \item \textbf{RQ2 (Interpretability).} Does increasing $R$ visibly reduce the attention-sink pattern in the CTC ViT encoder, and do per-character attention maps become better localized on the character they decode?
    \item \textbf{RQ3 (Architecture).} Which architectural choices are load-bearing for a CTC ViT HTR encoder (CNN stem, dual-head, RNN depth, patch layout)? Which are cosmetic?
    \item \textbf{RQ4 (Comparators).} How does the CTC ViT with registers compare, at equal training budget, against (i) a compact CNN--RNN baseline, (ii) a large ImageNet-pretrained ViT-B/16 fine-tuned on IAM, and (iii) a large in-domain-pretrained TrOCR-base fine-tuned on IAM?
\end{enumerate}

\subsection{Contributions}
\label{sec:contributions}

This report contributes:

\begin{enumerate}
    \item A \textbf{single-framework HTR pipeline} (\texttt{HTR-Pipeline}) that hosts four architecture families --- CNN--RNN, ViT-RGTS (with configurable $R$), pretrained TorchVision ViT-B/16, and pretrained TrOCR --- behind a shared training, evaluation, and attention-extraction interface, and an additional HTR-VT lane for the READ2016 external comparison.
    \item A \textbf{systematic register-token sweep} on IAM ($R \in \{0, 2, 4, 8, 16\}$, with and without stochastic weight averaging) that quantifies the CER-neutrality of registers for CTC HTR: the no-SWA CER range across $R$ is $5.74$--$5.98\%$ ($0.24$~pp), tightening to $5.60$--$5.70\%$ ($0.10$~pp) once the delivered SWA-averaged model is evaluated, and the seed-noise floor is $\pm 0.08$~pp over three seeds, all established from $150$-epoch runs.
    \item A \textbf{cross-dataset replication} on READ2016 that confirms the same near-flat CER-versus-$R$ pattern ($0.09$~pp range, $4.73$--$4.82\%$~CER), separating the finding from any IAM-specific artifact.
    \item An \textbf{attention-quality analysis} showing that register tokens do reduce the sink pattern qualitatively, that per-character CTC-column attention becomes measurably sharper with $R \geq 4$, and that a residual sink pattern remains --- traced to the fact that CTC-ViT patch tokens are simultaneously queries and keys of the sink attention, unlike the classification-ViT setting in which the CLS query and the patch sinks are structurally distinct.
    \item An \textbf{architecture ablation} that identifies the CNN stem and the strictly-left-to-right patch order it produces as the single non-negotiable design decision for a CTC ViT HTR encoder: removing the CNN stem (row-major $16{\times}16$ patches) collapses CER from $5.96\%$ to $73.36\%$, and the legacy v1 architecture with the same row-major layout independently collapses to $73.65\%$~CER --- two convergent failures on the same underlying cause.
    \item A \textbf{fine-tuning study} for the two pretrained comparators demonstrating that a two-group differential learning rate is required for ImageNet-pretrained ViT-B/16 ($15.09\%$~CER with two-group LR vs.\ $73.14\%$ with naive uniform LR) and that layer-wise learning-rate decay is required for TrOCR ($12.10\%$~CER with LLRD+SWA), while gradual unfreezing fails for both.
    \item An \textbf{honest domain-gap measurement} for synthetic pretraining: models trained purely on ${\sim}950$K synthetic lines reach ${\sim}0.5$--$0.8\%$~CER on synthetic-validation but $30$--$46\%$~CER on real IAM test, quantifying the render-to-real gap for the specific WordStylist-style synthetic corpus~\cite{nikolaidou2023wordstylist} used here.
    \item \textbf{Full reproducibility artifacts:} all $63$ training-run configurations, best-validation checkpoints, per-run \texttt{results.csv}/\texttt{training.log} pairs, attention weights for the register sweep runs, and the LaTeX result tables (\texttt{unified\_results.tex}, \texttt{register\_sweep.tex}, \texttt{ablations.tex}, \texttt{architecture\_comparison.tex}) shipped alongside the report.
\end{enumerate}

The rest of the report is organized as follows. Section~\ref{sec:related_work} surveys the HTR and register-token literature and positions this work against it. Section~\ref{sec:methodology} describes the datasets, the ViT-RGTS~v2 architecture, the register mechanism, and the training procedure. Section~\ref{sec:experiments} lays out the $63$-run experiment matrix and the evaluation protocol. Section~\ref{sec:results} reports the recognition, ablation, and attention-quality results, and includes a dedicated critical section on whether the Darcet~\etal\ claim holds in the CTC HTR setting. Section~\ref{sec:discussion} interprets the findings, and Section~\ref{sec:conclusion} closes with directions for extending the pipeline toward writer identification.
